\section{Configurations and Experimental Variables}
\label{sec:configurations}
Component configurations are managed via YAML and TOML templates (e.g., \texttt{sequencer.docker.yaml}, \texttt{sequencer.template.toml}). These templates are processed using \texttt{envsubst} to inject environment variables (like \texttt{SEQUENCER\_BIN}, \texttt{L1\_RPC\_URL}, and \texttt{BRIDGE\_ADDRESS}). It is crucial that default values are explicitly exported in the bash environment prior to substitution, as the \texttt{envsubst} implementation does not support bash fallback syntax.

\begin{table}[H]
\centering
\caption{Experiment Matrix: Stages 1 through 5}
\label{tab:experiment_matrix}
\small
\begin{tabularx}{\textwidth}{>{\hsize=0.5\hsize}X >{\hsize=1.2\hsize}X >{\hsize=1.3\hsize}X}
\toprule
\textbf{Stage} & \textbf{Focus Area} & \textbf{Key Independent Variables Swept} \\
\midrule
Stage 1 & Fixed Batching \& Latency-Cost Trade-off & Max Batch Size (10, 25, 50, 100, 250, 500, 1000) \\
\addlinespace
Stage 2 & Adaptive Batching Thresholds & Load Thresholds (Low, Medium, High) under varying offered TPS \\
\addlinespace
Stage 3 & Sequencer Scheduling Policies & Scheduling Algorithm (FCFS, FeePriority, TimeBoost, FairBFT) \\
\addlinespace
Stage 4 & Data Availability (DA) Modes & DA Mode (Calldata, Blob, Offchain), Blob Target Bytes \\
\addlinespace
Stage 5 & Prover Backend \& Cycle Analysis & Prover Mode (Mock vs RISC0), Target Batch Size (50, 100, 200, 500) \\
\bottomrule
\end{tabularx}
\end{table}

The metrics pipeline scripts (\texttt{data-tools/aggregate.py}) require these stages to be executed sequentially, relying on bash scripts inside \texttt{benchmark-suite/scripts/} (e.g., \texttt{run\_experiment.sh}). The pipeline joins metrics by batch ID across the \texttt{sequencer\_batch\_metrics.jsonl}, \texttt{executor\_batch\_metrics.jsonl}, and \texttt{submitter\_metrics.json} files to build a comprehensive lifecycle view of each transaction batch.